\paragraph{Estimating the global effect under interference.}
Randomisation-based inference under interference was developed around \emph{exposure mappings}, which reduce the full assignment vector to a coarser exposure; Horvitz--Thompson and H\'ajek estimators of exposure-specific means then follow from known exposure probabilities \citep{aronow2017estimating}. For social networks, graph cluster randomisation increases the probability that a unit's whole neighbourhood is treated, and the resulting HT estimators have a bias--variance trade-off governed by the clustering \citep{ugander2013graph,eckles2017design}; \citet{karrer2021network} describe such a system in production. \citet{leung2022rate} studies cluster designs and HT estimation of the global effect when interference decays with distance, and \citet{cortez2022staggered} show that staggered rollouts permit unbiased estimation of the total effect under low-order interactions without knowledge of the network. \citet{chen2023optimized} derive bias and variance of the HT estimator under one-hop interference and optimise the covariance of the assignment. Under unknown interference, standard estimators can still target a generalised estimand that averages over spillovers \citep{savje2021average}, and \citet{yu2022estimating} address estimation of the total effect when the network is unknown.

\paragraph{Regression-based estimators.}
The estimator studied here fits an outcome model on the assignment and graph-derived features from a single experiment and then predicts outcomes under all-treated and all-control assignments. \citet{chin2019regression} proposed this regression adjustment for Bernoulli designs, gave properties in linear models together with resampling-based inference, and evaluated it on a Facebook100 network. \citet{han2023model} generate and select graph- and treatment-based covariates, and construct block-bootstrap confidence intervals. These works show that regression adjustment can substantially reduce the bias of the difference in means when the model is adequate; they place more emphasis on the correctly specified and mildly misspecified linear regime than on characterising what determines the estimation error and where it breaks down, which is our focus.
